\documentclass[10pt,a4paper]{amsart}
\usepackage[margin=19mm]{geometry}
\usepackage{amsmath,amssymb,mathtools}
\usepackage[scr=rsfs,cal=euler]{mathalfa}
\usepackage{xfrac}
\usepackage[unicode,hidelinks,bookmarks=false]{hyperref}
\newcommand{\ie}{i.e.\ }
\newcommand{\cf}{cf.\ }
\newcommand{\resp}{respectively\ }
\newcommand{\Subsection}{\subsection}
\providecommand{\nobreakdash}{\nobreak\hbox{-}}
\setlength{\emergencystretch}{2em}
\multlinegap0pt
\usepackage{bm}
\usepackage{bbm}
\usepackage{dsfont}
\usepackage{xspace}
\usepackage{cancel}
\usepackage{mathtools}
\usepackage{amsthm}
\makeatletter
\@ifundefined{definition}{\newtheorem{definition}{Definition}}{}
\@ifundefined{theorem}{\newtheorem{theorem}{Theorem}}{}
\@ifundefined{proposition}{\newtheorem{proposition}[theorem]{Proposition}}{}
\makeatother
\def\R{{\mathbb {R}}}
\newcommand{\abs}[1]{\lvert#1\rvert}
\title[Nonexistence of global solutions to the Grushin heat equatio]{Nonexistence of global solutions to the Grushin heat equation with nonlocal and local nonlinearities}
\author{Ahmad Z. Fino, Arl\'ucio Viana}
\date{Source: arXiv:2510.09258v1 (CC BY 4.0)}
\begin{document}
\maketitle

\noindent\emph{Source and licence.} Nonexistence of global solutions to the Grushin heat equation with nonlocal and local nonlinearities. Version arXiv:2510.09258v1 (CC BY 4.0), distributed under the Creative Commons Attribution 4.0 International licence, as recorded in the arXiv licence metadata for this exact version and shown on the abstract page https://arxiv.org/abs/2510.09258v1. Full rights notice: licenses/arxiv-2510.09258-CC-BY-4.0.txt.\par\smallskip

\noindent\textit{Editorial scope.} Counted unit: the Definition whose source label is \texttt{weaksolution} in the article (source0.tex lines 382--393, source label \texttt{weaksolution}), delivered together with the article's own complete original proof of it (source0.tex lines 645--751, one \texttt{proof} environment of 107 lines). No mathematics is written, repaired or re-derived: statement and proof are the article's own text line for line. The proof is detached from the statement in the source: the article states the result at lines 382--393 and prints its proof at lines 645--751, and the proof environment's own header names the statement, so the association is the article's own. Required context retained uncounted: * (lines 133--140); IP (lines 238--252); Int3 (lines 258--267); Int4 (lines 264--266); Int6 (lines 297--299); IneqA (lines 550--553); eq2 (lines 621--629); T2 (lines 632--636). Every definition, display and label of those lines is retained unchanged. Omitted from this package and not redistributed: 1--132, 141--237, 253--257, 267--296, 300--381, 394--549, 554--620, 630--631, 637--644, 752--955. Nothing inside a retained excerpt is omitted or altered: every retained body is delivered exactly as the article's lines are written, so no omission receipt of any kind accompanies this package. Citations: the retained excerpts carry no citation command at all, so no bibliography entry of the article is transcribed and no reference is redistributed. Reference adaptation: none was needed and none was made; every reference inside the delivered unit resolves inside it, so no unresolved reference or citation is produced. Printed slips kept literal: no printed word, symbol, hypothesis or number of the counted statement or of its proof was altered. Layout and class: the article's own class is \texttt{amsart} with packages including amsmath, amsfonts, latexsym, amscd, amssymb, bbm, mathrsfs, hyperref, cleveref, xcolor, tikz, verbatim, setspace, ulem; the wrapper is the lane's pinned \texttt{amsart} editorial wrapper at 10pt on A4, which restates the theorem-like environments the retained text needs and adds the article's own macro definitions that the retained text needs (\texttt{\textbackslash R}, \texttt{\textbackslash abs}), with extra wrapper packages (bm, bbm, dsfont, xspace, cancel, mathtools). The delivered numbering is the unit's own. Assets: no retained span contains a figure environment, a graphics-inclusion command, a PDF-page inclusion, a table or a TikZ picture; no image file is copied into this package, the package declares no asset dependency and no image file is redistributed with it. Licence: version arXiv:2510.09258v1 of this article is distributed under the Creative Commons Attribution 4.0 International licence, as recorded in the arXiv licence metadata for this exact version and shown on the abstract page https://arxiv.org/abs/2510.09258v1. A full rights notice is delivered at licenses/arxiv-2510.09258-CC-BY-4.0.txt. Characters: every retained excerpt is pure ASCII, so the delivered engine, XeLaTeX with its default Latin Modern text font, reports no missing character.
\begin{equation}\label{*}
	\left\{\begin{array}{ll}
		\displaystyle {u_{t}-\Delta_{\mathcal{G}}
			u = k_1 \int_0^t(t-s)^{-\gamma}\abs{u}^{p_1-1}u(s)\,\mathrm{d}s} + k_2|u|^{p_2-1}u, &  (z,t)\in {\mathbb{R}}^{N+k}\times (0,\infty),\\
		\displaystyle{u(z,0)=  u_0(z),\qquad\qquad}&\displaystyle{z\in {\mathbb{R}}^{N+k},}
	\end{array}
	\right.
\end{equation}

\begin{proposition}
	Let $\alpha\in(0,1)$ and $-\infty<c<d<\infty$. The fractional integration by parts formula
	\begin{equation}\label{IP}
		\int_{c}^{d}f(t)D^\alpha_{c|t}g(t)\,\mathrm{dt} \;=\; \int_{c}^{d}
		g(t)D^\alpha_{t|d}f(t)\,\mathrm{dt},
	\end{equation}
	is satisfied for every $f\in I^\alpha_{t|d}(L^p(c,d))$, $g\in I^\alpha_{c|t}(L^q(c,d))$ such that $\frac{1}{p}+\frac{1}{\R^{N+k}}\leq 1+\alpha$, $p,q>1$, where
	\[
	I^\alpha_{c|t}(L^q(0,T)):=\left\{f= I^\alpha_{c|t}h,\,\, h\in L^q(c,d)\right\},
	\]
	and
	\[
	I^\alpha_{t|d}(L^p(c,d)):=\left\{f= I^\alpha_{t|d}h,\,\, h\in L^p(c,d)\right\}.
	\]
\end{proposition}

\begin{proposition}
	For $0<\alpha<1$, $-\infty<c<d<\infty$, we have the following identities
	\begin{equation}\label{Int3}
		D^\alpha_{c|t}I^\alpha_{c|t}\varphi(t)=\varphi(t),\,\hbox{a.e. $t\in(c,d)$}, \quad\hbox{for all}\,\,\varphi\in L^r(c,d),\,\, 1\leq r\leq\infty,
	\end{equation}
	and
	\begin{equation}\label{Int4}
		-\frac{d}{\mathrm{dt}} \, D^\alpha_{t|d}\varphi=D^{1+\alpha}_{t|d}\varphi,\quad \varphi\in AC^2[c,d].
	\end{equation}
\end{proposition}

	\begin{equation}\label{Int6}
		D_{t|T}^{m+\alpha}w_1(t)=\frac{\Gamma(\sigma+1)}{\Gamma(\sigma+1-m-\alpha)}T^{-(m+\alpha)}(1-t/T)^{\sigma-\alpha-m}.
	\end{equation}

	\begin{align}\label{IneqA}
		\Delta_{\mathcal{G}}\varphi(z,t) &=\ell(\ell-1)\varphi^\ell_2(y)\varphi^\ell_3(t)\varphi^{\ell-2}_1(x)|\nabla_x\varphi_1(x)|^2+\ell\varphi^\ell_2(y)\varphi^\ell_3(t)\varphi^{\ell-1}_1(x)\Delta_x\varphi_1(x) \\
		&{}\quad +\,\ell(\ell-1)|x|^2\varphi^\ell_1(x)\varphi^\ell_3(t)\varphi^{\ell-2}_2(y)|\nabla_y\varphi_2(y)|^2+\ell|x|^2\varphi^\ell_1(x)\varphi^\ell_3(t)\varphi^{\ell-1}_2(y)\Delta_y\varphi_2(y) .\nonumber
	\end{align}

\begin{equation}\label{eq2}
	\left\{\begin{array}{ll}
		\,\, \displaystyle {u_{t}-\Delta_{\mathcal{G}}
			u =\int_0^t(t-s)^{-\gamma}\abs{u}^{p-1}u(s)\,ds,} &\displaystyle {(z,t)\in {\mathbb{R}}^{N+k}\times (0,\infty),}\\
		{}\\
		\displaystyle{u(z,0)=  u_0(z),\qquad\qquad}&\displaystyle{z\in {\mathbb{R}}^{N+k},}
	\end{array}
	\right.
\end{equation} 

\begin{theorem}\label{T2} Assume $0< \gamma<1$, and let $u_0\in L^1(\R^{N+k})$ be such that $\displaystyle \int_{\R^{N+k}}u_0(z)\,\mathrm{dz}>0$. If
	$$p\leq1+\frac{2(2-\gamma)}{N+2k-2+2\gamma}:=p_0\quad\hbox{or}\quad
	p<\frac{1}{\gamma},$$
	then problem \eqref{eq2} admits no global nonnegative weak solution.
\end{theorem}

\begin{definition}[Weak solution]
	Let $u_0\in L_{\mathrm{loc}}^1(\R^{N+k})$, and $T>0$. We say that $u$ is a
	weak solution of the problem \eqref{*} if 
	$$I^\alpha_{0|t}(|u|^{p_1-1}u), |u|^{p_2-1}u \in L_{\mathrm{loc}}^1((0,T),L_{\mathrm{loc}}^1(\R^{N+k})) , $$
	and verifies the weak formulation
	\begin{eqnarray}\label{weaksolution}
		&&\int_{\R^{N+k}}u(z,\tau)\varphi(z,\tau)\,\mathrm{dz}-\int_{\R^{N+k}}u_0(z)\varphi(z,0)\,\mathrm{dz}\nonumber\\
		&&=k_1\Gamma(\alpha)\int_0^T\int_{\R^{N+k}}I^\alpha_{0|t}(|u|^{p_1-1}u)\varphi(z,t)\,\mathrm{dz}\,\mathrm{dt}+k_2\int_0^T\int_{\R^{N+k}}|u|^{p_2-1}u\varphi(z,t)\,\mathrm{dz}\,\mathrm{dt}\nonumber\\
		&&+\int_0^T\int_{\R^{N+k}}u(z,t)\Delta_{\mathcal{G}}\varphi(z,t)\,\mathrm{dz}\,\mathrm{dt}+\int_0^T\int_{\R^{N+k}}u(z,t)\varphi_t(z,t)\,\mathrm{dz}\,\mathrm{dt},
	\end{eqnarray}
	for all compactly supported $\varphi\in C^{2,1}(\R^{N+k}\times[0,T])$ such that $\varphi(T,\cdotp)=0$, where $\alpha:=1-\gamma\in(0,1)$.
\end{definition}

\begin{proof}[Proof of Theorem \ref{T2}]%label{subsecA}
	
	
	The proof is by contradiction. Suppose that $u$ is a global nonnegative weak
	solution to \eqref{eq2}, then,  \eqref{weaksolution} holds for all $T\gg1$, with $k_1=1$ and $k_2=0$. Let $R$ and $T$ be large parameters in $(0,\infty)$. Let us choose 
	$$\varphi(z,t)=D^\alpha_{t|T}\left(\widetilde{\varphi}(z,t)\right):=D^\alpha_{t|T}\left(\left(\varphi_1(x)\right)^\ell \left(\varphi_2(y)\right)^\ell\left(\varphi_3(t)\right)^\ell\right),\qquad z=(x,y)\in\mathbb{R}^{N+k},\,\,t\geq 0,$$
	with $\ell= 2p/(p-1)$,
	$$\varphi_1(x):=\Phi\left(\abs{x}/R^{1/2}\right),\quad \varphi_2(y):=\Phi\left(\abs{y}/R\right),\quad \varphi_3(t):=\left(1-t/T\right)_+,$$
	where $\Phi\in \mathcal{C}^\infty(\mathbb{R})$ is a smooth non-increasing function 
	satisfying $\mathbbm{1}_{(-\infty,1]} \leq \Phi\leq \mathbbm{1}_{(-\infty,2]}$.
	Thus
	\begin{eqnarray}\label{weak5}
		&&\int_{\Omega}u_0(x,y)\left(\varphi_1(x)\right)^\ell \left(\varphi_2(y)\right)^\ell D^\alpha_{t|T}(\varphi_3)(0)\,dx\,dy + \Gamma(\alpha) \int_{\Omega_T} I^\alpha_{0|t}(u^{p}) \left(D^{\alpha}_{t|T}\widetilde{\varphi}(z,t)\right)\,\mathrm{dz}\,\mathrm{dt}\nonumber\\
		&&=-\int_{\Omega_T} u(z,t) \left(\partial_tD^{\alpha}_{t|T}\widetilde{\varphi}(z,t)\right)\,\mathrm{dz}\,\mathrm{dt} -\int_{\Omega_T}u(z,t) \left(\Delta_{\mathcal{G}}D^{\alpha}_{t|T}\widetilde{\varphi}(z,t) \right)\,\mathrm{dz}\,\mathrm{dt}
	\end{eqnarray}
	where
	\[
	\Omega_T:=[0,T]\times\Omega\;\;\text{for}\;\; \Omega=\left\{ z=(x,y)\in{\mathbb{R}}^{N+k}\hspace{2
		mm};\hspace{2mm}|x|\leq 2 R^{1/2},\,|y|\leq 2 R\right\}.
	\]
	Furthermore, using \eqref{IP}, \eqref{Int3}, \eqref{Int4}, and \eqref{Int6} in  \eqref{weak5}, we obtain
	\begin{eqnarray}
		&&C\,T^{-\alpha}\int_{\Omega}u_0(x,y)\left(\varphi_1(x)\right)^\ell \left(\varphi_2(y)\right)^\ell \,dx\,dy+\Gamma(\alpha)\int_{\Omega_T} u^p(z,t) \widetilde{\varphi}(z,t)\,\mathrm{dz}\,\mathrm{dt} \nonumber\\
		&&=\int_{\Omega_T}u(z,t) \left(D^{1+\alpha}_{t|T}\widetilde{\varphi}(z,t)\right) \,\mathrm{dz}\,\mathrm{dt} -\int_{\Omega_T}u(z,t) \left(\Delta_{\mathcal{G}}D^{\alpha}_{t|T}\widetilde{\varphi}(z,t)\right)\,\mathrm{dz}\,\mathrm{dt} \nonumber\\
		&&=: J_1+J_2. \label{weak6}
	\end{eqnarray}
	\noindent In order to estimate $J_1$ and $J_2$, we introduce the term 
	$\widetilde{\varphi}^{1/p}\widetilde{\varphi}^{-1/p}$ and we use Young's inequality to obtain
	\begin{align}\label{J1}
		J_1\leqslant\frac{\Gamma(\alpha)}{4} \int_{\Omega_T}u^{p}(z,t)\widetilde{\varphi}(z,t)\,\mathrm{dz}\,\mathrm{dt}+C\int_{\Omega_T}\varphi_1^\ell(x)\varphi_2^\ell(y)\varphi_3^{-\ell p^\prime/p}(t)\left|D^{1+\alpha}_{t|T} \varphi^\ell_3(t)\right|^{p'}\mathrm{d}z\,\mathrm{d}t.
	\end{align}
	To estimate $J_2$, we apply identity \eqref{IneqA} and again the Young's inequality in the following way
	\begin{eqnarray}\label{J2}
		J_2&\leqslant&\frac{\Gamma(\alpha)}{4} \int_{\Omega_T}u^{p}(z,t)\widetilde{\varphi}(z,t)\,\mathrm{dz}\,\mathrm{dt}+C\int_{\Omega_T}\varphi_2^\ell(y)\varphi_3^{-\ell p^\prime/p}(t)\left|D^{\alpha}_{t|T} \varphi^\ell_3(t)\right|^{p'}\left|\nabla_x \varphi_1(x)\right|^{2p'}\mathrm{d}z\,\mathrm{d}t\notag\\
		&{}&+\,C\int_{\Omega_T}\varphi_2^\ell(y)\varphi_1^{p'}(x)\varphi_3^{-\ell p^\prime/p}(t)\left|D^{\alpha}_{t|T} \varphi^\ell_3(t)\right|^{p'}\left|\Delta_x \varphi_1(x)\right|^{p'}\mathrm{d}z\,\mathrm{d}t\nonumber\\
		&&+C\int_{\Omega_T}|x|^{2p'}\varphi_1^\ell(x)\varphi_3^{-\ell p^\prime/p}(t)\left|D^{\alpha}_{t|T} \varphi^\ell_3(t)\right|^{p'}\left|\nabla_y \varphi_2(y)\right|^{2p'}\mathrm{d}z\,\mathrm{d}t\notag\\
		&{}&+\,C\int_{\Omega_T}|x|^{2p'}\varphi_1^\ell(x)\varphi_2^{p'}(y)\varphi_3^{-\ell p^\prime/p}(t)\left|D^{\alpha}_{t|T} \varphi^\ell_3(t)\right|^{p'}\left|\Delta_y \varphi_2(y)\right|^{p'}\mathrm{d}z\,\mathrm{d}t.
	\end{eqnarray}
	Combining \eqref{weak6}, \eqref{J1}, and \eqref{J2}, we arrive at
	\begin{eqnarray}\label{weak7}
		&{}&C\,T^{-\alpha}\int_{\Omega}u_0(x,y)\left(\varphi_1(x)\right)^\ell \left(\varphi_2(y)\right)^\ell\,dx\,dy+\frac{\Gamma(\alpha)}{2}\int_{\Omega_T}u^{p}(z,t)\widetilde{\varphi}(z,t)\,\mathrm{dz}\,\mathrm{dt}\notag\\
		&{}&\leq C\int_{\Omega_T}\varphi_1^\ell(x)\varphi_2^\ell(y)\varphi_3^{-\ell p^\prime/p}(t)\left|D^{1+\alpha}_{t|T} \varphi^\ell_3(t)\right|^{p'}\mathrm{d}z\,\mathrm{d}t\notag\\&&\quad+\,C\int_{\Omega_T}\varphi_2^\ell(y)\varphi_3^{-\ell p^\prime/p}(t)\left|D^{\alpha}_{t|T} \varphi^\ell_3(t)\right|^{p'}\left|\nabla_x \varphi_1(x)\right|^{2p'}\mathrm{d}z\,\mathrm{d}t\notag\\
		&{}&\quad+\,\,C\int_{\Omega_T}\varphi_2^\ell(y)\varphi_1^{p'}(x)\varphi_3^{-\ell p^\prime/p}(t)\left|D^{\alpha}_{t|T} \varphi^\ell_3(t)\right|^{p'}\left|\Delta_x \varphi_1(x)\right|^{p'}\mathrm{d}z\,\mathrm{d}t\notag\\
		&&\quad+\,C\int_{\Omega_T}|x|^{2p'}\varphi_1^\ell(x)\varphi_3^{-\ell p^\prime/p}(t)\left|D^{\alpha}_{t|T} \varphi^\ell_3(t)\right|^{p'}\left|\nabla_y \varphi_2(y)\right|^{2p'}\mathrm{d}z\,\mathrm{d}t\notag\\
		&{}&\quad+\,\,C\int_{\Omega_T}|x|^{2p'}\varphi_1^\ell(x)\varphi_2^{p'}(y)\varphi_3^{-\ell p^\prime/p}(t)\left|D^{\alpha}_{t|T} \varphi^\ell_3(t)\right|^{p'}\left|\Delta_y \varphi_2(y)\right|^{p'}\mathrm{d}z\,\mathrm{d}t.
	\end{eqnarray}
	At this stage, we introduce the scaled variables 
	\begin{equation}\label{changevar}
		\tau_1=R^{-1/2}x,\quad \tau_2=R^{-1}y,\quad \tau_3=T^{-1}t,	
	\end{equation} 
	we get
	\begin{eqnarray}\label{weak8}
		&&C\,T^{-\alpha}\int_{\Omega}u_0(x,y)\left(\varphi_1(x)\right)^\ell \left(\varphi_2(y)\right)^\ell\,dx\,dy+\frac{\Gamma(\alpha)}{2}\int_{\Omega_T}u^{p}(z,t)\widetilde{\varphi}(z,t)\,\mathrm{dz}\,\mathrm{dt}\notag\\
		&&\leq C\,T^{1-(1+\alpha)p'}R^{\frac{N}{2}+k}+C\,T^{1-\alpha p'}R^{-p'+\frac{N}{2}+k}.
	\end{eqnarray}
	We have to distinguish three cases:\\
	
	\noindent $\bullet$ The case $p<p_0$. We divide the proof into two steps. In the first step, we take $R=T$. From \eqref{weak8}, and using the fact that $\varphi_1(x),\varphi_2(y)\leq 1$, we conclude that
	\begin{equation}\label{weak9}
		\int_{\Omega_T}(u(z,t))^{p}\widetilde{\varphi}(z,t)\,\mathrm{dz}\,\mathrm{dt}\leq C\,T^{-\delta}+C\,T^{-\alpha}\|u_0\|_{L^1}.
	\end{equation}
	with $\delta:=(1+\alpha)p'-\frac{N}{2}-k-1$. As $p<p_0$ is equivalent to $\delta>0$, we pass to the limit in \eqref{weak9} as $T\rightarrow\infty$, and using the monotone convergence theorem together with $u_0\in L^1(\mathbb{R}^{N+k})$, we get
	$$0\leq \int_0^\infty\int_{\R^{N+k}}(u(z,t))^{p}\,\mathrm{dz}\,\mathrm{dt}\leq 0.$$
	Therefore $u\equiv 0$ almost everywhere. In the second step, coming back to \eqref{weak6} combined with $u\equiv0$ a.e., we may arrive at
	$$C\,T^{-\alpha}\int_{\Omega}u_0(x,y)\left(\varphi_1(x)\right)^\ell \left(\varphi_2(y)\right)^\ell \,dx\,dy\leq 0,$$
	that is,
	$$\int_{\Omega}u_0(x,y)\left(\varphi_1(x)\right)^\ell \left(\varphi_2(y)\right)^\ell \,dx\,dy\leq 0,\quad\hbox{for all}\,\,T>0.$$
	Passing now to the limit when $T\rightarrow\infty$, and using Lebesgue dominated convergence theorem together with $u_0\in L^1(\mathbb{R}^{N+k})$, we get
	\[
	0<\int_{\R^{N+k}}u_0(x,y)\,dx\,dy=\lim_{T\rightarrow\infty}\int_{\Omega}u_0(x,y)\left(\varphi_{1,T}(x)\right)^\ell \left(\varphi_{2,T}(y)\right)^\ell\,dx\,dy\leq 0 ,
	\]
	a contradiction.\\
	
	\noindent $\bullet$ The case $p=p_0$. We can see from \eqref{weak9} with $T\rightarrow\infty$ and taking into account the fact that $p=p_0$ is equivalent to $\delta=0$, that
	\begin{equation}\label{regularityA}
		u\in L^p((0,\infty),L^p(\mathbb{R}^{N+k})).
	\end{equation}
	On the other hand, by repeating the same calculation performed after \eqref{weak6}, by using H\"older's inequality instead of Young's inequality in $J_2$, we arrive at
	\begin{eqnarray*}
		&{}&C\,T^{-\alpha}\int_{\Omega}u_0(x,y)\left(\varphi_1(x)\right)^\ell \left(\varphi_2(y)\right)^\ell \,dx\,dy+\Gamma(\alpha)\int_{\Omega_T}(u(z,t))^{p}\widetilde{\varphi}(z,t)\,\mathrm{dz}\,\mathrm{dt}\\
		&{}&\leq \frac{\Gamma(\alpha)}{4} \int_{\Omega_T}(u(z,t))^{p}\widetilde{\varphi}(z,t)\,\mathrm{dz}\,\mathrm{dt}+C\int_{\Omega_T}\varphi_1^\ell(x)\varphi_2^\ell(y)\varphi_3^{-\ell p^\prime/p}(t)\left|D^{1+\alpha}_{t|T} \varphi^\ell_3(t)\right|^{p'}\mathrm{d}z\,\mathrm{d}t\\
		&{}&+\left(\int_{\Sigma_T}(u(z,t))^{p}\widetilde{\varphi}(z,t)\,\mathrm{dz}\,\mathrm{dt}\right)^{1/p} \notag\\
		&&\times \left[ \left(\int_{\Omega_T}\varphi_2^\ell(y)\varphi_3^{-\ell p^\prime/p}(t)\left|D^{\alpha}_{t|T} \varphi^\ell_3(t)\right|^{p'}\left|\nabla_x \varphi_1(x)\right|^{2p'}\mathrm{d}z\,\mathrm{d}t\right)^{1/p'} \right.\notag\\
		&& \quad + \left(\int_{\Omega_T}\varphi_2^\ell(y)\varphi_1^{p'}(x)\varphi_3^{-\ell p^\prime/p}(t)\left|D^{\alpha}_{t|T} \varphi^\ell_3(t)\right|^{p'}\left|\Delta_x \varphi_1(x)\right|^{p'}\mathrm{d}z\,\mathrm{d}t\right)^{1/p'}\\
		&& \quad + \left(\int_{\Omega_T}|x|^{2p'}\varphi_1^\ell(x)\varphi_3^{-\ell p^\prime/p}(t)\left|D^{\alpha}_{t|T} \varphi^\ell_3(t)\right|^{p'}\left|\nabla_y \varphi_2(y)\right|^{2p'}\mathrm{d}z\,\mathrm{d}t\right)^{1/p'}\notag\\
		&& \quad + \left. \left(\int_{\Omega_T}|x|^{2p'}\varphi_1^\ell(x)\varphi_2^{p'}(y)\varphi_3^{-\ell p^\prime/p}(t)\left|D^{\alpha}_{t|T} \varphi^\ell_3(t)\right|^{p'}\left|\Delta_y \varphi_2(y)\right|^{p'}\mathrm{d}z\,\mathrm{d}t\right)^{1/p'} \right]
	\end{eqnarray*}
	where
	$$\Sigma_T:=[0,T]\times\Sigma\;\;\text{for}\;\; \Sigma=\left\{ z=(x,y)\in{\mathbb{R}}^{N+k}\hspace{2
		mm};\hspace{2mm}R^{1/2}\leq |x|\leq 2 R^{1/2},\,R\leq |y|\leq 2 R\right\}.$$
	Let $R=TK^{-1}$ with some constants $K\ge 1$ with $K<T$ such that $T$ and $K$ cannot go to infinity simultaneously. Using \eqref{changevar} and taking into account the fact that $p=p_c$, we get
	\begin{eqnarray*}
		&{}&\frac{3\Gamma(\alpha)}{4}\int_{\Omega_T}(u(z,t))^{p}\widetilde{\varphi}(z,t)\,\mathrm{dz}\,\mathrm{dt}\\
		&{}&\leq C\,T^{-\alpha}\|u_0\|_{L^1}+CK^{-\frac{N}{2}-k} +\,C\,K^{1-\frac{N}{2p'}-\frac{k}{p'}}\left(\int_{\Sigma_T}(u(z,t))^{p}\widetilde{\varphi}(z,t)\,\mathrm{dz}\,\mathrm{dt}\right)^{1/p}
	\end{eqnarray*}
	where we have used $\varphi_1(x),\varphi_2(y)\geq 0$. Passing to the limit as $T\rightarrow\infty$, and using \eqref{regularityA}, we conclude that
	$$\int_0^\infty\int_{\R^{N+k}}(u(z,t))^{p}\,\mathrm{dz}\,\mathrm{dt}\leq CK^{-\frac{N}{2}-k}.$$
	Therefore, taking a sufficiently large $K$ we obtain the desired contradiction similarly to the first case.\\
	
	\noindent $\bullet$ The case $p<1/\gamma$. We choose in this case $R<T$ such that $T$ and $R$ cannot go to infinity simultaneously. From \eqref{weak8}, and using the fact that $\varphi_1(x),\varphi_2(y)\leq 1$, we conclude that
	\begin{equation}\label{weak10}
		\int_{\Omega_T}(u(z,t))^{p}\widetilde{\varphi}_T(z,t)\,\mathrm{dz}\,\mathrm{dt}\leq C\,T^{1-(1+\alpha)p'}R^{\frac{N}{2}+k}+C\,T^{1-\alpha p'}R^{-p'+\frac{N}{2}+k}+C\,T^{-\alpha}\|u_0\|_{L^1}.
	\end{equation}
	Here we wrote $\widetilde{\varphi}_T$ to emphasize the dependence of $T$. As $p<1/\gamma\Longleftrightarrow 1-\alpha p'<0$, we pass to the limit in \eqref{weak10} as $T\rightarrow\infty$, and using $u_0\in L^1(\mathbb{R}^{N+k})$, we get
	$$0\leq \int_0^\infty\int_{\R^{N+k}}(u(z,t))^{p}\,\mathrm{dz}\,\mathrm{dt}\leq 0.$$
	Therefore $u\equiv 0$ almost everywhere, and so the desired contradiction can be obtained similarly to the first case.
\end{proof}

\end{document}
